\section{AMC 12 12A 2006}
        \begin{enumerate}
        	\item (\textit{AMC 12 12A 2006}) Sandwiches at Joe's Fast Food cost $3$ dollars each and sodas cost $2$ dollars each. How many dollars will it cost to purchase $5$ sandwiches and $8$ sodas?


 $\mathrm{(A) \ } 31\qquad \mathrm{(B) \ } 32\qquad \mathrm{(C) \ } 33\qquad \mathrm{(D) \ } 34\qquad \mathrm{(E) \ } 35$

 \textbf{Answer: }A



\textbf{Solution}

The $5$ sandwiches cost $5\cdot 3=15$ dollars. The $8$ sodas cost $8\cdot 2=16$ dollars. In total, the purchase costs $15+16=\boxed{\textbf{(A) }31}$ dollars.

 


	\item (\textit{AMC 12 12A 2006}) Define $x\otimes y=x^3-y$. What is $h\otimes (h\otimes h)$?


 $\mathrm{(A) \ } -h\qquad \mathrm{(B) \ } 0\qquad \mathrm{(C) \ } h\qquad \mathrm{(D) \ } 2h\qquad \mathrm{(E) \ }  h^3$

 \textbf{Answer: }C



\textbf{Solution 1}

By the definition of $\otimes$, we have $h\otimes h=h^{3}-h$. 

 Then, $h\otimes (h\otimes h)=h\otimes (h^{3}-h)=h^{3}-(h^{3}-h)=\boxed{\textbf{(C) }h}.$



	\item (\textit{AMC 12 12A 2006}) The ratio of Mary's age to Alice's age is $3:5$. Alice is $30$ years old. How old is Mary?


 $\mathrm{(A) \ } 15\qquad \mathrm{(B) \ } 18\qquad \mathrm{(C) \ } 20\qquad \mathrm{(D) \ } 24\qquad \mathrm{(E) \ }  50$

 \textbf{Answer: }B



\textbf{Solution 1}

Let $m$ be Mary's age. Then $\frac{m}{30}=\frac{3}{5}$. Solving for $m$, we obtain $m=\boxed{\textbf{(B) }18}.$




\textbf{Solution 2}

We can see this is a combined ratio of $8$, $(5+3)$. We can equalize by doing $30\div5=6$, and $6\cdot3=\boxed{\textbf{(B) }18}$. With the common ratio of $8$ and difference ratio of $6$, we see $6\cdot8=30+18$. Therefore, we can see our answer is correct.

 


	\item (\textit{AMC 12 12A 2006}) A digital watch displays hours and minutes with AM and PM. What is the largest possible sum of the digits in the display?


 $\mathrm{(A) \ } 17\qquad \mathrm{(B) \ } 19\qquad \mathrm{(C) \ } 21\qquad \mathrm{(D) \ } 22\qquad \mathrm{(E) \ }  23$

 \textbf{Answer: }E



\textbf{Solution 1}

From the \href{https://artofproblemsolving.com/wiki/index.php?title=Greedy\_algorithm}{greedy algorithm}, we have $9$ in the hours section and $59$ in the minutes section. $9+5+9=\boxed{\textbf{(E) }23}$




\textbf{Solution 2 (matrix)}

With a matrix, we can see$\begin{bmatrix} 1+2&9&6&3\\ 1+1&8&5&2\\ 1+0&7&4&1 \end{bmatrix}$The largest single digit sum we can get is $9$.For the minutes digits, we can combine the largest $2$ digits, which are $9,5 \Rightarrow 9+5=14$, and finally $14+9=\boxed{\textbf{(E) }23}$




\textbf{Solution 3}

We first note that since the watch displays time in AM and PM, the value for the hours section varies from $00-12$. Therefore, the maximum value of the digits for the hours is when the watch displays $09$, which gives us $0+9=9$.

 Next, we look at the value of the minutes section, which varies from $00-59$. Let this value be a number $ab$. We quickly find that the maximum value for $a$ and $b$ is respectively $5$ and $9$.

 Adding these up, we get $9+5+9=\boxed{\textbf{(E) }23}$.

 


	\item (\textit{AMC 12 12A 2006}) Doug and Dave shared a pizza with $8$ equally-sized slices. Doug wanted a plain pizza, but Dave wanted anchovies on half the pizza. The cost of a plain pizza was $8$ dollars, and there was an additional cost of $2$ dollars for putting anchovies on one half. Dave ate all the slices of anchovy pizza and one plain slice. Doug ate the remainder. Each paid for what he had eaten. How many more dollars did Dave pay than Doug?


 $\mathrm{(A) \ } 1\qquad \mathrm{(B) \ } 2\qquad \mathrm{(C) \ } 3\qquad \mathrm{(D) \ } 4\qquad \mathrm{(E) \ }  5$

 \textbf{Answer: }D



\textbf{Solution}

Dave and Doug paid $8+2=10$ dollars in total.  Doug paid for three slices of plain pizza, which cost $\frac{3}{8}\cdot 8=3$.  Dave paid $10-3=7$ dollars.  Dave paid $7-3=\boxed{\textbf{(D) }4}$ more dollars than Doug.

 


	\item (\textit{AMC 12 12A 2006}) The $8\times 18$ rectangle $ABCD$ is cut into two congruent hexagons, as shown, in such a way that the two hexagons can be repositioned without overlap to form a square. What is $y$?

                    \begin{figure}[H]
                        \centering
                        \includegraphics[width=0.5\textwidth]{Problem/AMC_12_12A_2006_6_Problem_Figure_1.png}
                    \end{figure}
                    
$\mathrm{(A) \ } 6\qquad \mathrm{(B) \ } 7\qquad \mathrm{(C) \ } 8\qquad \mathrm{(D) \ } 9\qquad \mathrm{(E) \ }  10$

 \textbf{Answer: }A



\textbf{Solution 1}

Since the two \href{https://artofproblemsolving.com/wiki/index.php?title=Hexagon}{hexagons} are going to be repositioned to form a \href{https://artofproblemsolving.com/wiki/index.php?title=Square\_(geometry)}{ square} without overlap, the \href{https://artofproblemsolving.com/wiki/index.php?title=Area}{area} will remain the same. The rectangle's area is $18\cdot8=144$. This means the square will have four sides of length 12. The only way to do this is shown below.


 
                    \begin{figure}[H]
                        \centering
                        \includegraphics[width=0.5\textwidth]{Solution/AMC_12_12A_2006_6_Solution_Figure_1.png}
                    \end{figure}
                    

 As you can see from the diagram, the \href{https://artofproblemsolving.com/wiki/index.php?title=Line\_segment}{line segment} denoted as $y$ is half the length of the side of the square, which leads to $y = \frac{12}{2} = \boxed{\textbf{(A) }6}$.

 



\textbf{Solution 2  (Shortcut)}

As solution 1 says, the two hexagons are going to be repositioned to form a square without overlap. Thus we create this square out of the original rectangle.

 

                    \begin{figure}[H]
                        \centering
                        \includegraphics[width=0.5\textwidth]{Solution/AMC_12_12A_2006_6_Solution_Figure_2.png}
                    \end{figure}
                    

 As you can see from the diagram, the length $y$ fits into the previously blank side, so we know that it is equal to $y$.

 
                    \begin{figure}[H]
                        \centering
                        \includegraphics[width=0.5\textwidth]{Solution/AMC_12_12A_2006_6_Solution_Figure_3.png}
                    \end{figure}
                    

 
From there we can say $3y = 18$ so $y = \frac{18}{3} = \boxed{\textbf{(A) }6}$.

 



\textbf{Solution 3 (Cheap)}

Because the two hexagons are congruent, we know that the perpendicular line to $A$ is half of $BC$, or $4$. Next, we plug the answer choices in to see which one works. Trying $A$, we get the area of one hexagon is $72$ , as desired, so the answer is $\boxed{\textbf{(A) }6}$ .

 


	\item (\textit{AMC 12 12A 2006}) Mary is $20\%$ older than Sally, and Sally is $40\%$ younger than Danielle. The sum of their ages is $23.2$ years. How old will Mary be on her next birthday?


 $\mathrm{(A) \ } 7\qquad \mathrm{(B) \ } 8\qquad \mathrm{(C) \ } 9\qquad \mathrm{(D) \ } 10\qquad \mathrm{(E) \ }  11$

 \textbf{Answer: }B



\textbf{Solution}

Let $m$ be Mary's age, let $s$ be Sally's age, and let $d$ be Danielle's age.  

 We have $s=.6d$, so $m=1.2s=1.2(.6d)=.72d$.  

 The sum of their ages is $m+s+d=.72d+.6d+d=2.32d$.  Therefore, $2.32d=23.2$, and $d=10$.

 The smallest positive integer more than $7.2$ is $8$, so Mary will be $\boxed {\textbf {(B) } 8}$ on her next birthday.

 


	\item (\textit{AMC 12 12A 2006}) How many sets of two or more consecutive positive integers have a sum of $15$?


 $\mathrm{(A) \ } 1\qquad \mathrm{(B) \ } 2\qquad \mathrm{(C) \ } 3\qquad \mathrm{(D) \ } 4\qquad \mathrm{(E) \ }  5$

 \textbf{Answer: }C



\textbf{Solution 1}

Notice that if the consecutive positive integers have a sum of $15$, then their average (which could be a fraction) must be a divisor of $15$. If the number of integers in the list is odd, then the average must be either $1, 3,$ or $5$, and $1$ is clearly not possible. The other two possibilities both work:

 
            \begin{itemize}
        \item $1 + 2 + 3 + 4 + 5 = 15$
        \item $4 + 5 + 6 = 15$
            \end{itemize}
            If the number of integers in the list is even, then the average will have a $\frac{1}{2}$. The only possibility is $\frac{15}{2}$, from which we get:

 
            \begin{itemize}
        \item $15 = 7 + 8$
            \end{itemize}
            Thus, the correct answer is $\boxed{\textbf{(C) }3}.$

 
Question: (RealityWrites) Is it possible that the answer is $4$, because $0+1+2+3+4+5$ should technically count, right?

 Answer: (IMGROOT2) It isn't possible because the question asks for positive integers, and this means that negative integers or zero aren't allowed.

 Note to readers: make sure to always read the problem VERY carefully before attempting; it could mean the difference of making the cutoff.

 



\textbf{Solution 2}

Any set will form a arithmetic progression with the first term say $a$. Since the numbers are consecutive the common difference $d = 1$.

 The sum of the AP has to be 15. So,

 
 \[S_n = \frac{n}{2} \cdot (2a + (n-1)d)\]
 \[S_n = \frac{n}{2} \cdot (2a + (n-1)1)\]
 \[15 = \frac{n}{2} \cdot (2a + n - 1)\]
 \[2an + n^2 - n = 30\]
 \[n^2 + n(2a - 1) - 30 = 0\]

 Now we need to find the number of possible sets of values of a, n which satisfy this equation. Now $a$ cannot be 15 as we need 2 terms. So a can only be less the 15.

 Trying all the values of a from 1 to 14 we observe that $a = 4$, $a = 7$ and $a = 1$ provide the only real solutions to the above equation.The three possibilites of a and n are.

 
 \[(a,n) = (4,3),(7, 2),(1, 6)\]

 The above values are obtained by solving the following equations obtained by substituting the above mentioned values of a into $n^2 + n(2a - 1) - 30 = 0$,

 
 \[n^2 + 7n - 30 = 0\]
 \[n^2 + 13n - 30 = 0\]
 \[n^2 - n - 30 = 0\]

 Since there are 3 possibilities the answer is $\boxed{\textbf{(C) }3}.$




\textbf{Solution 3}

We want some consecutive numbers that add up. We can use casework to solve this. We know that any set of 6 consecutive, positive integers cannot be greater than 5. Therefore the max amount of elements is 5.

 \textbf{Case 1: There are 5 numbers} -We know that 15/5 = 3, so if we put 3 as our middle number and expand evenly, we get {1,2,3,4,5}

 \textbf{Case 2: There are 4 numbers} -By a bit of bashing we find that there are no solutions

 \textbf{Case 3: There are 3 numbers} -We know that 15/3 = 5, so if we put 5 as our middle number and expand evenly, we get {4,5,6}

 \textbf{Case 4: There are 2 numbers} -We'll say x is the lower number. Thus x + (x + 1) = 15. Solving for x, x = 7. Thus getting {7,8}

 Adding up all the cases we get 1 + 0 + 1 + 1 = $\boxed{\textbf{(C) }3}.$



	\item (\textit{AMC 12 12A 2006}) Oscar buys $13$ pencils and $3$ erasers for $\textdollar 1.00$. A pencil costs more than an eraser, and both items cost a whole number of cents. What is the total cost, in cents, of one pencil and one eraser?


 $\mathrm{(A) \ } 10\qquad \mathrm{(B) \ } 12\qquad \mathrm{(C) \ } 15\qquad \mathrm{(D) \ } 18\qquad \mathrm{(E) \ }  20$

 \textbf{Answer: }A



\textbf{Solution 1}

Let the price of a pencil be $p$ and an eraser $e$.  Then $13p + 3e = 100$ with $p > e > 0$.  Since $p$ and $e$ are \href{https://artofproblemsolving.com/wiki/index.php?title=Positive\_integer}{positive integers}, we must have $e \geq 1$ and $p \geq 2$.

 Considering the \href{https://artofproblemsolving.com/wiki/index.php?title=Equation}{equation} $13p + 3e = 100$ \href{https://artofproblemsolving.com/wiki/index.php?title=Modulo}{modulo} 3 (that is, comparing the \href{https://artofproblemsolving.com/wiki/index.php?title=Remainder}{remainders} when both sides are divided by 3) we have $p + 0e \equiv 1 \pmod 3$ so $p$ leaves a remainder of 1 on division by 3.

 Since $p \geq 2$, possible values for $p$ are 4, 7, 10 ....

 Since 13 pencils cost less than 100 cents, $13p < 100$.  $13 \times 10 = 130$ is too high, so $p$ must be 4 or 7.

 If $p = 4$ then $13p = 52$ and so $3e = 48$ giving $e = 16$.  This contradicts the pencil being more expensive.  The only remaining value for $p$ is 7; then the 13 pencils cost $7 \times 13= 91$ cents and so the 3 erasers together cost 9 cents and each eraser costs $\frac{9}{3} = 3$ cents.

 Thus one pencil plus one eraser cost $7 + 3 = 10$ cents, which is answer choice $\mathrm{(A) \ }$.

 



\textbf{Solution 2}

Since we know that the values of pencils and erasers are both whole numbers, and that there are $3$ erasers and $13$ pencils, we can choose numbers to be the value for the pencil, and if the remainder to $1.00$ is divisible by 3, we know our answer is correct. There are 7 numbers to check, although we can eliminate $1, 2, 3, 4,$ and $5$ mentally, since those numbers would probably have the cost of erasers higher than the cost of pencils (6 would too, but it is faster to just test 6). Trying 6, $13 \cdot 6 = 72$. 28 is not divisible by 3, so we know that this number is not correct. Moving on to 7, $13 \cdot 7 = 91$. We know that 9 is a multiple of 3, so this value works! 

 Since the total cost of erasers is 9 cents, we know the individual eraser is $\frac{9}{3} = 3$. $3 + 7 = 10$ so our answer is answer choice $\mathrm{(A) \ }$



	\item (\textit{AMC 12 12A 2006}) For how many real values of $x$ is $\sqrt{120-\sqrt{x}}$ an integer?


 $\mathrm{(A) \ } 3\qquad \mathrm{(B) \ } 6\qquad \mathrm{(C) \ } 9\qquad \mathrm{(D) \ } 10\qquad \mathrm{(E) \ }  11$

 \textbf{Answer: }E



\textbf{Solution}

For $\sqrt{120-\sqrt{x}}$ to be an integer, $120-\sqrt{x}$ must be a perfect square. 

 Since $\sqrt{x}$ can't be negative, $120-\sqrt{x} \leq 120$. 

 The perfect squares that are less than or equal to $120$ are $\{0,1,4,9,16,25,36,49,64,81,100\}$, so there are $11$ values for $120-\sqrt{x}$. 

 Since every value of $120-\sqrt{x}$ gives one and only one possible value for $x$, the number of values of $x$ is $\boxed{\textbf{(E) }11}$.

 


	\item (\textit{AMC 12 12A 2006}) Which of the following describes the graph of the equation $(x+y)^2=x^2+y^2$?


 $\mathrm{(A)}\ \text{the empty set}\qquad\mathrm{(B)}\ \text{one point}\qquad\mathrm{(C)}\ \text{two lines}\qquad\mathrm{(D)}\ \text{a circle}\qquad\mathrm{(E)}\ \text{the entire plane}$

 \textbf{Answer: }C



\textbf{Solution}


 \begin{align*}(x+y)^2&=x^2+y^2\\ x^2 + 2xy + y^2 &= x^2 + y^2\\ 2xy &= 0\end{align*}
Either $x = 0$ or $y = 0$. The \href{https://artofproblemsolving.com/wiki/index.php?title=Union}{union} of them is 2 lines, and thus the answer is $\mathrm{(C)}$.

 
                    \begin{figure}[H]
                        \centering
                        \includegraphics[width=0.5\textwidth]{Solution/AMC_12_12A_2006_11_Solution_Figure_1.png}
                    \end{figure}
                    



	\item (\textit{AMC 12 12A 2006}) A number of linked rings, each 1 cm thick, are hanging on a peg. The top ring has an outside diameter of 20 cm. The outside diameter of each of the outer rings is 1 cm less than that of the ring above it. The bottom ring has an outside diameter of 3 cm. What is the distance, in cm, from the top of the top ring to the bottom of the bottom ring?

                    \begin{figure}[H]
                        \centering
                        \includegraphics[width=0.5\textwidth]{Problem/AMC_12_12A_2006_12_Problem_Figure_1.png}
                    \end{figure}
                    
$\mathrm{(A) \ } 171\qquad \mathrm{(B) \ } 173\qquad \mathrm{(C) \ } 182\qquad \mathrm{(D) \ } 188\qquad \mathrm{(E) \ }  210$

 \textbf{Answer: }B



\textbf{Solution 1}

The inside diameters of the rings are the positive integers from $1$ to $18$.  The total distance needed is the sum of these values plus $2$ for the top of the first ring and the bottom of the last ring. Using the formula for the sum of an \href{https://artofproblemsolving.com/wiki/index.php?title=Arithmetic\_series}{arithmetic series}, the answer is $\frac{18 \cdot 19}{2} + 2 = \boxed{\textbf{(B) }173}$.

 



\textbf{Solution 2}

Alternatively, the sum of the consecutive integers from 3 to 20 is $\frac{1}{2}(18)(3+20) = 207$.  However, the 17 \href{https://artofproblemsolving.com/wiki/index.php?title=Intersection}{intersections} between the rings must be subtracted, and we also get $207 - 2(17) = \boxed{\textbf{(B) }173}$.

 



\textbf{Solution 3 (Brute force)}

We can just add all the numbers up. But, we have to account for the intersections, so the final expression is $3+\sum_{k=2}^{18} k$, which is $\boxed{\textbf{(B) }173}$. ~\href{https://artofproblemsolving.com/wiki/index.php?title=Shunyipanda&action=edit&redlink=1}{shunyipanda}

 


	\item (\textit{AMC 12 12A 2006}) The vertices of a $3-4-5$ right triangle are the centers of three mutually externally tangent circles, as shown. What is the sum of the areas of the three circles?

                    \begin{figure}[H]
                        \centering
                        \includegraphics[width=0.5\textwidth]{Problem/AMC_12_12A_2006_13_Problem_Figure_1.png}
                    \end{figure}
                    
$\mathrm{(A) \ } 12\pi\qquad \mathrm{(B) \ } \frac{25\pi}{2}\qquad \mathrm{(C) \ } 13\pi\qquad \mathrm{(D) \ } \frac{27\pi}{2}\qquad \mathrm{(E) \ }  14\pi$

 \textbf{Answer: }E



\textbf{Solution}

Let the radius of the smallest circle be $r_A$, the radius of the second largest circle be $r_B$, and the radius of the largest circle be $r_C$. 
 \[r_A + r_B = 3\]
 \[r_A + r_C = 4\]
 \[r_ B + r_C = 5\]

 Adding up all these equations and then dividing both sides by 2, we get,

 
 \[r_A + r_B + r_C = 6\]

 Then, we get $r_A = 1$, $r_B = 2$, and $r_C = 3$ Then we get $1^2 \pi + 2^2 \pi + 3^2 \pi = 14 \pi \iff\mathrm{(E)}$



	\item (\textit{AMC 12 12A 2006}) Two farmers agree that pigs are worth $300$ dollars and that goats are worth $210$ dollars. When one farmer owes the other money, he pays the debt in pigs or goats, with "change" received in the form of goats or pigs as necessary. (For example, a $390$ dollar debt could be paid with two pigs, with one goat received in change.) What is the amount of the smallest positive debt that can be resolved in this way?


 $\mathrm{(A) \ } \textdollar5 \qquad \mathrm{(B) \ } \textdollar 10 \qquad \mathrm{(C) \ } \textdollar 30 \qquad \mathrm{(D) \ } \textdollar 90 \qquad \mathrm{(E) \ }  \textdollar 210$

 \textbf{Answer: }C



\textbf{Solution 1 (Diophantine Equation)}

The problem can be restated as an equation of the form $300p + 210g = x$, where $p$ is the number of pigs, $g$ is the number of goats, and $x$ is the positive debt. The problem asks us to find the lowest \textit{x} possible. \textit{$p$} and \textit{$g$} must be integers, which makes the equation a \href{https://artofproblemsolving.com/wiki/index.php?title=Diophantine\_equation}{Diophantine equation}.  \href{https://artofproblemsolving.com/wiki/index.php?title=Bezout\%27s\_Lemma}{Bezout's Lemma} tells us that the smallest $c$ for the Diophantine equation $am + bn = c$ to have solutions is when $c$ is the GCD (\href{https://artofproblemsolving.com/wiki/index.php?title=Greatest\_common\_divisor}{greatest common divisor}) of $a$ and $b$. Therefore, the answer is $gcd(300,210)=\boxed{\textbf{(C) }30}.$




\textbf{Solution 2 (Divisibility Analysis)}

Alternatively, note that $300p + 210g = 30(10p + 7g)$ is divisible by $30$ no matter what $p$ and $g$ are, so our answer must be divisible by $30$. Since we want the smallest integer, we can suppose that the answer is $30$ and go on from there. Note that three goats minus two pigs gives us $630 - 600 = 30$ exactly.  Since our supposition can be achieved, the answer is $\boxed{\textbf{(C) }30}$.

 



\textbf{Solution 3 (Simplifying the Problem)}

Let us simplify this problem.  Dividing by $30$, we get a pig to be:  $\frac{300}{30} =  10$, and a goat to be $\frac{210}{30}= 7$. It becomes evident that if you exchange $5$ pigs for $7$ goats, we get the smallest positive difference - $5\cdot 10 - 7\cdot 7 = 50-49 = 1$, since we can't make a non-integer with a linear combination of integers.  Since we originally divided by $30$, we need to multiply again, thus getting the answer  $1\cdot 30 = \boxed{\textbf{(C) }30}$.

 


	\item (\textit{AMC 12 12A 2006}) Suppose $\cos x=0$ and $\cos (x+z)=1/2$. What is the smallest possible positive value of $z$?


 $\mathrm{(A) \ } \frac{\pi}{6}\qquad \mathrm{(B) \ } \frac{\pi}{3}\qquad \mathrm{(C) \ } \frac{\pi}{2}\qquad \mathrm{(D) \ } \frac{5\pi}{6}\qquad \mathrm{(E) \ }  \frac{7\pi}{6}$

 \textbf{Answer: }A



\textbf{Solution}


            \begin{itemize}
        \item For $\cos x = 0$, x must be in the form of $\frac{\pi}{2} + \pi n$, where $n$ denotes any \href{https://artofproblemsolving.com/wiki/index.php?title=Integer}{integer}.
        \item For $\cos (x+z) = 1 / 2$, $x + z = \frac{\pi}{3} +2\pi n, \frac{5\pi}{3} + 2\pi n$.
            \end{itemize}
            The smallest possible value of $z$ will be that of $\frac{5\pi}{3} - \frac{3\pi}{2} = \frac{\pi}{6} \Rightarrow \mathrm{(A)}$.

 



\textbf{Solution 2}

Notice that $\cos (x+z) = \cos x \cos z - \sin x \sin z$, we will obtain $\cos (x+z) = - \sin x \sin z$ after plugging $\cos x = 0$. Knowing that $\cos (x+z) = \frac{1}{2}$, we have $\sin x \sin z = - \frac{1}{2}$. 

 Now we can try plugging the answer choices back in to see if each works: 

 When $z = \frac{\pi}{6}$, $\sin z = \frac{1}{2}$ and $\sin x = - 1$. Since $\sin x$ is obtainable when $x = \frac{3\pi}{2}$, $z = \frac{\pi}{6}$ works. 

 The question asks for the smallest positive value of $z$, and $\frac{\pi}{6}$ is the smallest among the choices. Thus, the answer is $\boxed{\textbf{(A) } \frac{\pi}{6}}$



	\item (\textit{AMC 12 12A 2006}) Circles with centers $A$ and $B$ have radii $3$ and $8$, respectively. A common internal tangent intersects the circles at $C$ and $D$, respectively. Lines $AB$ and $CD$ intersect at $E$, and $AE=5$. What is $CD$?

                    \begin{figure}[H]
                        \centering
                        \includegraphics[width=0.5\textwidth]{Problem/AMC_12_12A_2006_16_Problem_Figure_1.png}
                    \end{figure}
                    
$\mathrm{(A)}\ 13\qquad\mathrm{(B)}\ \frac{44}{3}\qquad\mathrm{(C)}\ \sqrt{221}\qquad\mathrm{(D)}\ \sqrt{255}\qquad\mathrm{(E)}\ \frac{55}{3}$

 \textbf{Answer: }B



\textbf{Solution}


                    \begin{figure}[H]
                        \centering
                        \includegraphics[width=0.5\textwidth]{Solution/AMC_12_12A_2006_16_Solution_Figure_1.png}
                    \end{figure}
                    
$\triangle ADE \sim \triangle BCE$ because
 \[\angle AED \cong \angle BEC \text{ (vertical angles)}\]
 \[\angle ADE \cong \angle BCE \text{ (tangent perpendicularity)}\] 

 By the \href{https://artofproblemsolving.com/wiki/index.php?title=Pythagorean\_Theorem}{Pythagorean Theorem}, line segment $DE = 4$.  The sides are proportional, so 
 \[\frac{DE}{AD} = \frac{CE}{BC} \Rightarrow \frac{4}{3} = \frac{CE}{8}\Rightarrow CE = \frac{32}{3}\]

 And hence:
 \[CD = CE + DE = 4 + \frac{32}{3} = \boxed{\textbf{(B) }\frac{44}{3}}\]



	\item (\textit{AMC 12 12A 2006}) Square $ABCD$ has side length $s$, a circle centered at $E$ has radius $r$, and $r$ and $s$ are both rational. The circle passes through $D$, and $D$ lies on $\overline{BE}$. Point $F$ lies on the circle, on the same side of $\overline{BE}$ as $A$. Segment $AF$ is tangent to the circle, and $AF=\sqrt{9+5\sqrt{2}}$. What is $r/s$?

                    \begin{figure}[H]
                        \centering
                        \includegraphics[width=0.5\textwidth]{Problem/AMC_12_12A_2006_17_Problem_Figure_1.png}
                    \end{figure}
                    
$\mathrm{(A) \ } \frac{1}{2}\qquad \mathrm{(B) \ } \frac{5}{9}\qquad \mathrm{(C) \ } \frac{3}{5}\qquad \mathrm{(D) \ } \frac{5}{3}\qquad \mathrm{(E) \ }  \frac{9}{5}$

 \textbf{Answer: }B



\textbf{Solution 1}

One possibility is to use the \href{https://artofproblemsolving.com/wiki/index.php?title=Coordinate\_plane}{coordinate plane}, setting $B$ at the origin. Point $A$ will be $(0,s)$ and $E$ will be $\left(s + \frac{r}{\sqrt{2}},\ s + \frac{r}{\sqrt{2}}\right)$ since $B, D$, and $E$ are \href{https://artofproblemsolving.com/wiki/index.php?title=Collinear}{collinear} and contain a diagonal of $ABCD$. The \href{https://artofproblemsolving.com/wiki/index.php?title=Pythagorean\_theorem}{Pythagorean theorem} results in

 
 \[AF^2 + EF^2 = AE^2\]

 \[r^2  + \left(\sqrt{9 + 5\sqrt{2}}\right)^2  = \left(\left(s + \frac{r}{\sqrt{2}}\right) - 0\right)^2 + \left(\left(s + \frac{r}{\sqrt{2}}\right) - s\right)^2\]

 \[r^2  + 9 + 5\sqrt{2} = s^2 + rs\sqrt{2} + \frac{r^2}{2} + \frac{r^2}{2}\]

 \[9 + 5\sqrt{2} = s^2 + rs\sqrt{2}\]

 This implies that $rs = 5$ and $s^2 = 9$; dividing gives us $\frac{r}{s} = \frac{5}{9} \Rightarrow B$.

 



\textbf{Solution 2}

First note that angle $\angle AFE$ is right since $\overline{AF}$ is tangent to the circle. Using the Pythagorean Theorem on $\triangle AFE$, then, we see
 \[AE^2 = 9 + 5\sqrt{2} + r^2.\]

 But it can also be seen that $\angle BDA = 45^\circ$. Therefore, since $D$ lies on $\overline{BE}$, $\angle ADE = 135^\circ$. Using the Law of Cosines on $\triangle ADE$, we see

 
 \[AE^2 = s^2 + r^2 - 2sr\cos(135^\circ)\]
 \[AE^2 = s^2 + r^2 - 2sr\left(-\frac{1}{\sqrt{2}}\right)\]
 \[AE^2 = s^2 + r^2 + \sqrt{2}sr\]
 \[9 + 5\sqrt{2} + r^2 = s^2 + r^2 + \sqrt{2}sr\]
 \[9 + 5\sqrt{2} = s^2 + \sqrt{2}sr.\]

 Thus, since $r$ and $s$ are rational, $s^2 = 9$ and $sr = 5$. So $s = 3$, $r = \frac{5}{3}$, and $\frac{r}{s} = \frac{5}{9}$.

 



\textbf{Solution 3}

(Similar to Solution 1)First, draw line AE and mark a point Z that is equidistant from E and D so that $\angle{DZE} = 90^\circ$ and that line $\overline{AZ}$ includes point D. Since DE is equal to the radius $r$, $DZ=EZ=\frac{r}{\sqrt2}=\frac{r\sqrt2}{2}.$

 Note that triangles $\triangle AFE$ and $\triangle AZE$ share the same hypotenuse $(AE)$, meaning that
 \[AZ^2+EZ^2=AF^2+EF^2\]Plugging in our values we have:
 \[\left(s+\frac{r\sqrt{2}}{2}\right)^2+\left(\frac{r\sqrt{2}}{2}\right)^2=\left(\sqrt{9+5\sqrt{2}}\right)^2+r^2\]
 \[s^2+sr\sqrt{2}+\frac{r^2}{2}+\frac{r^2}{2}=9+5\sqrt{2}+r^2\]
 \[s^2+sr\sqrt{2}=9+5\sqrt{2}\]By logic $s=3$ and $sr=5 \implies r=5/3.$

 Therefore, $\frac{\frac{5}{3}}{3}=\frac{5}{9}=\boxed{B}$




\textbf{Solution 4 - Alcumus}

Let $B=(0,0)$, $C=(s,0)$, $A=(0,s)$, $D=(s,s)$, and $E=\left(s+\frac{r}{\sqrt{2}},s+\frac{r}{\sqrt{2}} \right)$. Apply the Pythagorean Theorem to $\triangle AFE$ to obtain 
 \[r^2+\left(9+5\sqrt{2}\right)=\left(s+\frac{r}{\sqrt{2}}\right)^2+\left(\frac{r}{\sqrt{2}}\right)^2,\] from which $9+5\sqrt{2}=s^2+rs\sqrt{2}$. Because $r$ and $s$ are rational, it follows that $s^2=9$ and $rs=5$, so $r/s = \boxed{5/9}$.

 OR

 Extend $\overline{AD}$ past $D$ to meet the circle at $G \ne D$. Because $E$ is collinear with $B$ and $D$, $\angle EDG = 45^\circ.$ Also, $ED = EG,$ which implies $\angle EGD = 45^\circ$, so $\triangle EDG$ is an isosceles right triangle. Thus $DG = r\sqrt{2}$. By the Power of a Point Theorem,
 \begin{align*} 9+5\sqrt{2} &= AF^2 \\ &= AD\cdot AG\\ & = AD\cdot \left(AD+DG\right) \\ &= s\left(s+r\sqrt{2}\right) \\ &= s^2+rs\sqrt{2}.\end{align*}
 As in the first solution, we conclude that $r/s=\boxed{5/9}$.

 



\textbf{Solution 5 - Power of a Point}

Extend $\overline{AD}$ from D to hit $\odot E$ at G. Let the foot of the altitude of $\triangle DEG$ be $M$.

 
The problem states that point $D$ is on $\overline{BE}$, so $\angle ADB \cong \angle EDM$. As $\angle DAB \cong \angle OMD$, we have $\triangle DAB \cong \triangle DMO$ by AA similarity. Note that $\triangle ADB$ is isosceles right; this implies that $\triangle DMO$ is as well. We have that $DM = \frac{ r \sqrt{2} }{2}$. As $\triangle DOG$ is clearly isosceles, $\overline{OM}$ is also a median, so $DM = GM$, yielding that $DG = r\sqrt{2}$.

 
From here, we consider the power of point $A$, observing that $(AF)^2 = (AD)(AG)$, or $9 + 5\sqrt{2} = s(s + r\sqrt{2} )$. Expanding, we get $s^2 = 9$ and $sr = 5$, so $s= 3$, $r = \frac{5}{3}$, and $\frac{r}{s} = \frac{ \frac{5}{3} }{3} = \frac{5}{9}$.

 



\textbf{Solution 6 - Answer Choices}

(Last minute solution) We roughly measure the distance of $r$ and the distance of $y$. Since $r$ is clearly less than $s$, we can eliminate answer choices (D) and (E). Next, if we compare the distances, $r$ seems to be just a little more than half of $s$, thus eliminating answer choice (A). $r$ is only a little bit bigger than half of $s$, so we can reasonably assume that their ratio is less than $\frac{3}{5}$. That leaves us with answer choice $\boxed{C}$ , or $\frac{5}{9}$.

 


	\item (\textit{AMC 12 12A 2006}) The function $f$ has the property that for each real number $x$ in its domain, $1/x$ is also in its domain and 


 $f(x)+f\left(\frac{1}{x}\right)=x$


 What is the largest set of real numbers that can be in the domain of $f$?


 $\mathrm{(A) \ } \{x|x\ne 0\}\qquad \mathrm{(B) \ } \{x|x<0\}\qquad \mathrm{(C) \ } \{x|x>0\}\qquad \mathrm{(D) \ } \{x|x\ne -1\;$ $\mathrm{and}\; x\ne 0\;\mathrm{and}\; x\ne 1\}\qquad \mathrm{(E) \ }  \{-1,1\}$

 \textbf{Answer: }E



\textbf{Solution}

Quickly verifying by plugging in values verifies that $-1$ and $1$ are in the domain.

 $f(x)+f\left(\frac{1}{x}\right)=x$

 Plugging in $\frac{1}{x}$ into the \href{https://artofproblemsolving.com/wiki/index.php?title=Function}{function}: 

 $f\left(\frac{1}{x}\right)+f\left(\frac{1}{\frac{1}{x}}\right)=\frac{1}{x}$
$f\left(\frac{1}{x}\right)+ f(x)= \frac{1}{x}$

 Since $f(x) + f\left(\frac{1}{x}\right)$ cannot have two values: 

 $x = \frac{1}{x}$
$x^2 = 1$
$x=\pm 1$

 Therefore, the largest \href{https://artofproblemsolving.com/wiki/index.php?title=Set}{set} of \href{https://artofproblemsolving.com/wiki/index.php?title=Real\_number}{real numbers} that can be in the \href{https://artofproblemsolving.com/wiki/index.php?title=Domain}{domain} of $f$ is $\{-1,1\} \Rightarrow E$




\textbf{Solution 2}

We know that $f(x) + f \left(\frac{1}{x}\right) = x.$ Plugging in $x = \frac{1}{x}$ we get 
 \[f \left(\frac{1}{x}\right) + f \left(\frac{1}{\frac{1}{x}}\right) = \frac{1}{x}\] 
 \[f \left(\frac{1}{x}\right) + f(x) = \frac{1}{x}.\]

 Also notice 
 \[f \left(\frac{1}{x}\right) + f(x) = x\] by the commutative property(this is the same as the equation given in the problem. We are just rearranging). So we can set $\frac{1}{x} = x$ which gives us $x = \pm 1$ which is answer option $\boxed{\mathrm{(E) \ }  \{-1,1\}}.$



	\item (\textit{AMC 12 12A 2006}) Circles with centers $(2,4)$ and $(14,9)$ have radii $4$ and $9$, respectively. The equation of a common external tangent to the circles can be written in the form $y=mx+b$ with $m>0$. What is $b$?


 
                    \begin{figure}[H]
                        \centering
                        \includegraphics[width=0.5\textwidth]{Problem/AMC_12_12A_2006_19_Problem_Figure_1.png}
                    \end{figure}
                    
$\mathrm{(A) \ } \frac{908}{119}\qquad \mathrm{(B) \ } \frac{909}{119}\qquad \mathrm{(C) \ } \frac{130}{17}\qquad \mathrm{(D) \ } \frac{911}{119}\qquad \mathrm{(E) \ }  \frac{912}{119}$

 \textbf{Answer: }E



\textbf{Solution}

Let $L_1$ be the line that goes through $(2,4)$ and $(14,9)$, and let $L_2$ be the line $y=mx+b$. If we let $\theta$ be the measure of the acute angle formed by $L_1$ and the x-axis, then $\tan\theta=\frac{5}{12}$. $L_1$ clearly bisects the angle formed by $L_2$ and the x-axis, so $m=\tan{2\theta}=\frac{2\tan\theta}{1-\tan^2{\theta}}=\frac{120}{119}$. We also know that $L_1$ and $L_2$ intersect at a point on the x-axis. The equation of $L_1$ is $y=\frac{5}{12}x+\frac{19}{6}$, so the coordinate of this point is $\left(-\frac{38}{5},0\right)$. Hence the equation of $L_2$ is $y=\frac{120}{119}x+\frac{912}{119}$, so $b=\frac{912}{119}$, and our answer choice is $\boxed{\mathrm{E}}$.

 


	\item (\textit{AMC 12 12A 2006}) A bug starts at one vertex of a cube and moves along the edges of the cube according to the following rule. At each vertex the bug will choose to travel along one of the three edges emanating from that vertex. Each edge has equal probability of being chosen, and all choices are independent. What is the probability that after seven moves the bug will have visited every vertex exactly once?


 $\mathrm{(A) \ } \frac{1}{2187}\qquad \mathrm{(B) \ } \frac{1}{729}\qquad \mathrm{(C) \ } \frac{2}{243}\qquad \mathrm{(D) \ } \frac{1}{81}\qquad \mathrm{(E) \ }  \frac{5}{243}$

 \textbf{Answer: }C



\textbf{Solution 1}

Call this cube $ABCDEFGH$, with $A$ being the starting point.

 Because in $7$ moves the bug has to visit the other vertices, the bug cannot revisit any vertex.

 Therefore, starting at $A$, the bug has a $\frac{3}{3}$ chance of finding a good path to the next vertex, and call it $B$.

 Then, the bug has a $\frac{2}{3}$ chance of reaching a new vertex next. Call this $C$. $A, B,$ and $C$ are always on the same plane.

 Now, we split cases.

 Case $1$: The bug goes to the vertex $E$ opposite $A$ on the space diagonal with probability $\frac{1}{3}$.

 Then, the bug has to visit $D$ on the plane of $ABC$ last, as there is no way in and out from $D$.

 Therefore, there is only $1$ way out of $81$ to get to $D$ last.

 Therefore, there is a $\frac{1}{243} \cdot \frac{6}{9} = \frac{6}{2187}$ chance of finding a good path in this case.

 Case $2$: The bug goes to vertex $D$ on plane $ABC$ with a chance of $\frac{1}{3}$.

 The bug then has only $1$ way to go to a point $E$ on the opposite face, therefore having a $\frac{1}{3}$ probability.

 Then, the bug has a choice of two vertices on the face opposite to $ABCD$.

 This results in a $\frac{2}{3}$ probability of finding a good path to a point $F$.

 Then, there is only $1$ way out of $9$ to visit both other vertices on that face in $2$ moves.

 Multiply the probabilities for this case to get $\frac{2}{243} \cdot \frac{6}{9} = \frac{12}{2187}$.

 Add the probabilities of these two cases together to get $\frac{18}{2187} = \boxed{\textbf{(C) }\frac{2}{243}}.$




\textbf{Solution 2}


                    \begin{figure}[H]
                        \centering
                        \includegraphics[width=0.5\textwidth]{Solution/AMC_12_12A_2006_20_Solution_Figure_1.png}
                    \end{figure}
                    

 Let us count the good paths. The bug starts at an arbitrary \href{https://artofproblemsolving.com/wiki/index.php?title=Vertex}{vertex}, moves to a neighboring vertex ($3$ ways), and then to a new neighbor ($2$ more ways).  So, without loss of generality, let the cube have vertices $ABCDEFGH$ such that $ABCD$ and $EFGH$ are two opposite faces with $A$ above $E$ and $B$ above $F$.  The bug starts at $A$ and moves first to $B$, then to $C$.

 From this point, there are two cases. 

 Case $1$: the bug moves to $D$.  From $D$, there is only one good move available, to $H$.  From $H$, there are two ways to finish the trip, either by going $H \to G \to F \to E$ or $H \to E \to F \to G$.  So there are $2$ good paths in this case.

 Case $2$: the bug moves to $G$.

 Case $2a$: the bug moves $G \to H$.  In this case, there are $0$ good paths because it will not be possible to visit both $D$ and $E$ without double-visiting some vertex.

 Case $2b$: the bug moves $G \to F$.  There is $1$ good path in this case, $F \to E \to H \to D$.

 Thus, there are a total of $3\cdot 3 \cdot 2 = 18$ good paths.  There were $3^7 = 2187$ possible paths the bug could have taken, so the probability of a random path is good is $\frac{18}{2187} = \boxed{\textbf{(C) }\frac2{243}}$.

 



\textbf{Solution 3 (answer choices)}

As in Solution 1, the bug can move from its arbitrary starting vertex to a neighboring vertex in $3$ ways. After this, the bug can move to a new neighbor in $2$ ways (it cannot return to the first vertex). The total number of paths (as stated above) is $3^7$ or $2187$. Therefore, the probability of the bug following a good path is equal to $\frac{6x}{2187}$ for some positive integer $x$. The only answer choice which can be expressed in this form is $\boxed{\textbf{(C) }\frac2{243}}$.

 



\textbf{Solution 4}

Pick 1 vertex to start from. Notice that on all the odd moves the bug will move to 1 of 4 other vertices. Notice that on all even moves the bug must go to 1 of 3 other vertices (4 including the starting vertex). Thus, the bug will end on one of the odd vertices. There are 3 choices for the bug's first move from the starting vertex, and 2 choices for the bug's second move (it can't go back to the previous vertices). Then, notice that the bug has 3 choices of an odd vertex to end on not including the vertex the bug went to on move 1. Now pick 1 of the 3 ending vertices to end on. Notice that there are only 6 valid paths from start to finish (1 valid path after moves 1 and 2) to that point because all the moves are forced after the first 2 moves  (try it out). Thus, there are $3 \cdot 2 \cdot 3 = 18$ valid paths the bug can take. There are $3^7 = 2187$total paths, so the answer is 18/2187 = $\boxed{\textbf{(C) }\frac{2}{243}}$.

 



\textbf{Solution 5 (classical probability)}

Let the cube have vertices $ABCDEFGH$ such that $ABCD$ and $EFGH$ are two opposite faces with $A$ above $E$ and $B$ above $F$.  The bug starts at $A$ and every step has 3 choices so that there're $3^7 = 2187$ possible paths. We can split the case into odd and even steps. After an odd number of steps, the possible points that can be reached are $B$,$D$,$G$,$E$. Then, for even-numbered steps, the remaining points can be reached. Since he has taken seven steps, the end point must be one of the points in $B$,$D$,$G$,$E$. At this point, we classify and discuss the movement trajectory of the bug. 

 Case 1: It moves to another plane after completing one plane. If it starts on the upper plane, there are two possibilities for the first step, $B$ or $D$. The second and third steps are restricted to the upper plane, and the fourth step will move to the nearby point on another plane. These three steps have only one possibility each, so there is no need for classification. Then, when it reaches the lower plane, the fifth step has two possibilities, thus limiting the sixth and seventh steps as well. Proceeding in this way, the bug can also move from left to right or from front to back. Therefore, there are $2 \cdot 2 \cdot 3 = 12$ situations. 

 Case 2: Opposite to case 1, it can directly pass through to another point G on the other diagonal of the cube within the first three steps. There are three choices for the first step, two choices for the second step, and only one choice for the third step. Therefore, there are $3 \cdot 2 = 6$ situations for the first three steps. The subsequent steps will be restricted to a single path (try it yourself). 

 In summary, there are 12 + 6 = 18 situations in total. Therefore, the final probability is 18/2187 = $\boxed{\textbf{(C) }\frac2{243}}$.

 



\textbf{Solution 6}

We can tell that the total number of paths will be $3^6$ as the first and last vertex of the cube will be forced. We can now let the bottom left vertex of the front face of the cube be $A$ and, without loss of generality, assume the bug starts here. Now we can list out the possibilities for the bug to touch each vertex once after $7$ moves starting at $A$; we find $6$ favorable outcomes of the $729$, so we find the answer to be $\frac{6}{729} = \boxed{\textbf{(C) }\frac{2}{243}}$.

 


	\item (\textit{AMC 12 12A 2006}) Let 


 $S_1=\{(x,y)|\log_{10}(1+x^2+y^2)\le 1+\log_{10}(x+y)\}$


 and 


 $S_2=\{(x,y)|\log_{10}(2+x^2+y^2)\le 2+\log_{10}(x+y)\}$.


 What is the ratio of the area of $S_2$ to the area of $S_1$?


 $\mathrm{(A) \ } 98\qquad \mathrm{(B) \ } 99\qquad \mathrm{(C) \ } 100\qquad \mathrm{(D) \ } 101\qquad \mathrm{(E) \ }  102$

 \textbf{Answer: }E



\textbf{Solution}

Looking at the constraints of $S_1$:

 $x+y > 0$
$\log_{10}(1+x^2+y^2)\le 1+\log_{10}(x+y)$
$\log_{10}(1+x^2+y^2)\le \log_{10} 10 +\log_{10}(x+y)$
$\log_{10}(1+x^2+y^2)\le \log_{10}(10x+10y)$
$1+x^2+y^2 \le 10x+10y$
$x^2-10x+y^2-10y \le -1$
$x^2-10x+25+y^2-10y+25 \le 49$
$(x-5)^2 + (y-5)^2 \le (7)^2$
$S_1$ is a circle with a radius of $7$. So, the area of $S_1$ is $49\pi$.

 Looking at the constraints of $S_2$:

 $x+y > 0$
$\log_{10}(2+x^2+y^2)\le 2+\log_{10}(x+y)$
$\log_{10}(2+x^2+y^2)\le \log_{10} 100 +\log_{10}(x+y)$
$\log_{10}(2+x^2+y^2)\le \log_{10}(100x+100y)$
$2+x^2+y^2 \le 100x+100y$
$x^2-100x+y^2-100y \le -2$
$x^2-100x+2500+y^2-100y+2500 \le 4998$
$(x-50)^2 + (y-50)^2 \le (7\sqrt{102})^2$
$S_2$ is a circle with a radius of $7\sqrt{102}$. So, the area of $S_2$ is $4998\pi$.

 So the desired ratio is $\frac{4998\pi}{49\pi} = 102 \Rightarrow \boxed{E}$.

 


	\item (\textit{AMC 12 12A 2006}) A circle of radius $r$ is concentric with and outside a regular hexagon of side length $2$. The probability that three entire sides of hexagon are visible from a randomly chosen point on the circle is $1/2$. What is $r$?


 $\mathrm{(A) \ } 2\sqrt{2}+2\sqrt{3}\qquad \mathrm{(B) \ } 3\sqrt{3}+\sqrt{2}\qquad \mathrm{(C) \ } 2\sqrt{6}+\sqrt{3} \qquad \mathrm{(D) \ } 3\sqrt{2}+\sqrt{6}\qquad \mathrm{(E) \ }  6\sqrt{2}-\sqrt{3}$

 \textbf{Answer: }D



\textbf{Solution}

Project any two non-adjacent and non-opposite sides of the \href{https://artofproblemsolving.com/wiki/index.php?title=Hexagon}{hexagon} to the \href{https://artofproblemsolving.com/wiki/index.php?title=Circle}{circle}; the \href{https://artofproblemsolving.com/wiki/index.php?title=Arc}{arc} between the two points formed is the location where all three sides of the hexagon can be fully viewed. Since there are six such pairs of sides, there are six arcs. The probability of choosing a point is $1 / 2$, or if the total arc degree measures add up to $\frac{1}{2} \cdot 360^{\circ} = 180^{\circ}$. Each arc must equal $\frac{180^{\circ}}{6} = 30^\mathrm{\circ}$.

 
                    \begin{figure}[H]
                        \centering
                        \includegraphics[width=0.5\textwidth]{Solution/AMC_12_12A_2006_22_Solution_Figure_1.png}
                    \end{figure}
                    

 Call the center $O$, and the two endpoints of the arc $A$ and $B$, so $\angle AOB = 30^{\circ}$. Let $P$ be the intersections of the projections of the sides of the hexagon corresponding to $\overline{AB}$. Notice that $\triangle APO$ is an \href{https://artofproblemsolving.com/wiki/index.php?title=Isosceles\_triangle}{isosceles triangle}: $\angle AOP = 15^{\circ}$ and $\angle OAP = OAB - 60^{\circ} = \frac{180^{\circ}-30^{\circ}}{2} - 60^{\circ} = 15^{\circ}$. Since $OA$ is a radius and $OP$ can be found in terms of a side of the hexagon, we are almost done.

 If we draw the altitude of $APO$ from $P$, then we get a \href{https://artofproblemsolving.com/wiki/index.php?title=Right\_triangle}{right triangle}. Using simple trigonometry, $\cos 15^{\circ} = \frac{\frac{r}{2}}{2\sqrt{3}} = \frac{r}{4\sqrt{3}}$.

 Since $\cos 15^{\circ} = \cos (45^{\circ} - 30^{\circ}) = \frac{\sqrt{6} + \sqrt{2}}{4}$, we get $r = \left(\frac{\sqrt{6} + \sqrt{2}}{4}\right) \cdot 4\sqrt{3} = 3\sqrt{2} + \sqrt{6} \Rightarrow \boxed{D}$.

 


	\item (\textit{AMC 12 12A 2006}) Given a finite sequence $S=(a_1,a_2,\ldots ,a_n)$ of $n$ real numbers, let $A(S)$ be the sequence 


 $\left(\frac{a_1+a_2}{2},\frac{a_2+a_3}{2},\ldots ,\frac{a_{n-1}+a_n}{2}\right)$


 of $n-1$ real numbers. Define $A^1(S)=A(S)$ and, for each integer $m$, $2\le m\le n-1$, define $A^m(S)=A(A^{m-1}(S))$. Suppose $x>0$, and let $S=(1,x,x^2,\ldots ,x^{100})$. If $A^{100}(S)=(1/2^{50})$, then what is $x$?


 $\mathrm{(A) \ } 1-\frac{\sqrt{2}}{2}\qquad \mathrm{(B) \ } \sqrt{2}-1\qquad \mathrm{(C) \ } \frac{1}{2}\qquad \mathrm{(D) \ } 2-\sqrt{2}\qquad \mathrm{(E) \ }  \frac{\sqrt{2}}{2}$

 \textbf{Answer: }B



\textbf{Solution 1}


 \[A^1(S)=\left(\frac{1+x}{2},\frac{x+x^2}{2},...,\frac{x^{99}+x^{100}}{2}\right)\]
 \[A^2(S)=\left(\frac{1+2x+x^2}{2^2},\frac{x+2x^2+x^3}{2^2},...,\frac{x^{98}+2x^{99}+x^{100}}{2^2}\right)\]
 \[\Rightarrow A^2(S)=\left(\frac{(x+1)^2}{2^2},\frac{x(x+1)^2}{2^2},...,\frac{x^{98}(x+1)^2}{2^2}\right)\]

 In general, $A^n(S)=\left(\frac{(x+1)^n}{2^n},\frac{x(x+1)^n}{2^n},...,\frac{x^{100-n}(x+1)^n}{2^n}\right)$ such that $A^n(s)$ has $101-n$ terms. Specifically, $A^{100}(S)=\frac{(x+1)^{100}}{2^{100}}$ To find x, we need only solve the equation $\frac{(x+1)^{100}}{2^{100}}=\frac{1}{2^{50}}$. Algebra yields $x=\sqrt{2}-1$.

 



\textbf{Solution 2}

For every sequence $S=\left(a_1,a_2,\dots, a_n\right)$ of at least three terms,

 
 \[A^2(S)=\left(\frac{a_1+2a_2+a_3}{4},\frac{a_2+2a_3+a_4}{4},\dots,\frac{a_{n-2}+2a_{n-1}+a_n}{4}\right).\]Thus for $m = 1\text{ and }2$, the coefficients of the terms in the numerator of $A^m(S)$ are the binomial coefficients $\binom{m}{0},\binom{m}{1},\dots,\binom{m}{m}$, and the denominator is $2^m$. Because $\binom{m}{r}+\binom{m}{r+1}=\binom{m+1}{r+1}$ for all integers $r\geq 0$, the coefficients of the terms in the numerators of $A^{m+1}(S)$ are $\binom{m+1}{0},\binom{m+1}{1},\ldots,\binom{m+1}{m+1}$ for $2\leq m\leq n-2$. The definition implies that the denominator of each term in $A^{m+1}(S)$ is $2^{m+1}$. For the given sequence, the sole term in $A^{100}(S)$ is 
 \[\frac{1}{2^{100}} \sum_{m=0}^{100} \binom{100}{m}a_{m+1} = \frac{1}{2^{100}} \sum_{m=0}^{100} \binom{100}{m}x^m = \frac{1}{2^{100}}(x+1)^{100}.\]Therefore, 
 \[\left(\frac{1}{2^{50}}\right)=A^{100}(S)=\left(\frac{(1+x)^{100}}{2^{100}}\right),\]so $(1+x)^{100}=2^{50}$, and because $x>0$, we have $x=\boxed{\sqrt{2}-1}$.- Alcumus

 


	\item (\textit{AMC 12 12A 2006}) The expression 


 $(x+y+z)^{2006}+(x-y-z)^{2006}$


 is simplified by expanding it and combining like terms. How many terms are in the simplified expression?


 $\mathrm{(A) \ } 6018\qquad \mathrm{(B) \ } 671,676\qquad \mathrm{(C) \ } 1,007,514\qquad \mathrm{(D) \ } 1,008,016\qquad \mathrm{(E) \ }  2,015,028$

 \textbf{Answer: }D



\textbf{Solution 1}

By the \href{https://artofproblemsolving.com/wiki/index.php?title=Multinomial\_Theorem}{Multinomial Theorem}, the summands can be written as

 
 \[\sum_{a+b+c=2006}{\frac{2006!}{a!b!c!}x^ay^bz^c}\]

 and

 
 \[\sum_{a+b+c=2006}{\frac{2006!}{a!b!c!}x^a(-y)^b(-z)^c},\]

 respectively. Since the coefficients of like terms are the same in each expression, each like term either cancel one another out or the coefficient doubles. In each expansion there are:

 
 \[{2006+2\choose 2} = 2015028\]

 terms without cancellation. For any term in the second expansion to be negative, the parity of the exponents of $y$ and $z$ must be opposite. Now we find a pattern:

 if the exponent of $y$ is $1$, the exponent of $z$ can be all even integers up to $2004$, so there are $1003$ terms.

 if the exponent of $y$ is $3$, the exponent of $z$ can go up to $2002$, so there are $1002$ terms.

 $\vdots$

 if the exponent of $y$ is $2005$, then $z$ can only be 0, so there is $1$ term.

 If we add them up, we get $\frac{1003\cdot1004}{2}$ terms. However, we can switch the exponents of $y$ and $z$ and these terms will still have a negative sign. So there are a total of $1003\cdot1004$ negative terms.

 By subtracting this number from 2015028, we obtain $\boxed{\mathrm{D}}$ or $1,008,016$ as our answer.

 



\textbf{Solution 2}

Alternatively, we can use a \href{https://artofproblemsolving.com/wiki/index.php?title=Generating\_function}{generating function} to solve this problem. The goal is to find the generating function for the number of unique terms in the simplified expression (in terms of $k$). In other words, we want to find $f(x)$ where the coefficient of $x^k$ equals the number of unique terms in $(x+y+z)^k + (x-y-z)^k$.

 
First, we note that all unique terms in the expression have the form, $Cx^ay^bz^c$, where $a+b+c=k$ and $C$ is some constant. Therefore, the generating function for the MAXIMUM number of unique terms possible in the simplified expression of $(x+y+z)^k + (x-y-z)^k$ is
 \[(1+x+x^2+x^3\cdots)^3 = \frac{1}{(1-x)^3}\]

 
Secondly, we note that a certain number of terms of the form, $Cx^ay^bz^c$, do not appear in the simplified version of our expression because those terms cancel. Specifically, we observe that terms cancel when $1 \equiv b+c\text{ (mod }2\text{)}$ because every unique term is of the form:
 \[\binom{k}{a,b,c}x^ay^bz^c+(-1)^{b+c}\binom{k}{a,b,c}x^ay^bz^c\]for all possible $a,b,c$.

 
Since the generating function for the maximum number of unique terms is already known, it is logical that we want to find the generating function for the number of terms that cancel, also in terms of $k$. With some thought, we see that this desired generating function is the following:
 \[2(x+x^3+x^5\cdots)(1+x^2+x^4\cdots)(1+x+x^2+x^3\cdots) = \frac{2x}{(1-x)^3(1+x)^2}\]

 
Now, we want to subtract the latter from the former in order to get the generating function for the number of unique terms in $(x+y+z)^k + (x-y-z)^k$, our initial goal:
 \[\frac{1}{(1-x)^3}-\frac{2x}{(1-x)^3(1+x)^2} = \frac{x^2+1}{(1-x)^3(1+x)^2}\]which equals
 \[(x^2+1)(1+x+x^2\cdots)^3(1-x+x^2-x^3\cdots)^2\]

 
The coefficient of $x^{2006}$ of the above expression equals
 \[\sum_{a=0}^{2006}\binom{2+a}{2}\binom{1+2006-a}{1}(-1)^a + \sum_{a=0}^{2004}\binom{2+a}{2}\binom{1+2004-a}{1}(-1)^a\]

 
Evaluating the expression, we get $1008016$, as expected.

 



\textbf{Solution 3}

Define $P$ such that $P=y+z$. Then the expression in the problem becomes:$(x+P)^{2006}+(x-P)^{2006}$. 

 Expanding this using binomial theorem gives $(x^n+Px^{n-1}+...+P^{n-1}x+P^n)+(x^n-Px^{n-1}+...-P^{n-1}x+P^n)$, where $n=2006$ (we may omit the coefficients, as we are seeking for the number of terms, not the terms themselves). 

 Simplifying gives: $2(x^n+x^{n-2}P^2+...+x^2P^{n-2}+P^n)$. Note that two terms that come out of different powers of $P$ cannot combine and simplify, as their exponent of $x$ will differ. Therefore, we simply add the number of terms produced from each addend. By the Binomial Theorem, $P^k=(y+z)^k$ will have $k+1$ terms, so the answer is $1+3+5+...+2007=1004^2=1,008,016 \implies \boxed{D}$.

 



\textbf{Solution 4}

Using stars and bars we know that $(x+y+z)^{2006}$ has ${2006+2\choose 2}$ or $2015028$ different terms. Now we need to find out how many of these terms are canceled out by $(x-y-z)^{2006}$. We know that for any term(let's say $x^{a}(-y)^{b}(-z)^{c}$) where $a+b+c=2006$ of the expansion of  $(x-y-z)^{2006}$ is only going to cancel out with the corresponding term $x^{a}y^{b}z^{c}$ if only $b$ is odd and $c$ is even or $b$ is even and $c$ is odd. Now let's do some casework to see how many terms fit this criteria:

 Case 1:    $a$ is even  

 Now we know that $a$ is even and $a+b+c=2006$. Thus $b+c$ is also even or both $b$ and $c$ are odd or both $b$ and $c$ are even. This case clearly fails the above criteria and there are 0 possible solutions.

 Case 2:    $a$ is odd

 Now we know that $a$ is odd and $a+b+c=2006$. Thus $b+c$ is odd and $b$ is odd and $c$ is even or $b$ is even and $c$ is odd. All terms that have $a$ being odd work.

 We now need to figure out how many of the terms have $a$ as a odd number. We know that $a$ can be equal to any number between 0 and 2006. There are 1003 odd numbers between this range and 2007 total numbers. Thus $\frac{1003}{2007}$ of the $2015028$ terms will cancel out and  $\frac{1004}{2007}$ of the terms will work. Thus there are $(\frac{1004}{2007})(2015028)$ terms. This number comes out to be $1,008,016$ $\implies \boxed{D}$(Author: David Camacho)

 



\textbf{Solution 5}

Noticing how $y$ and $z$ are negative in the second part of the expression, let $x=a$ and $y+z = b$. Then we get 
 \[(a+b)^{2006} + (a-b)^{2006}\]We know that the terms that don't cancel out have even powers of $a$ and $b$. Our expansion will be in the form of
 \[2a^{2006} + x_1a^{2004}b^{2} + x_2a^{2002}b^{4} + \cdots + 2b^{2006}\]Note that $b^n = (x+y)^n$ has $n+1$ terms. Furthermore, the current expression is irreducible as each term has a different $x$ power. Thus, when we write $a$ and $b$ back to their original terms, we will have $1+3+5+ \cdots + 2007 = 1004^2 = 1,008,016 \implies \boxed{D}$

 -smartguy888

 



\textbf{Solution 6}

Only consider $(x-y-z)^{2006}$ for now. If a term there is subtracted, it would be canceled out since there is a corresponding positive term in $(x+y+z)^{2006}$. So the only terms that will appear are the terms with positive coefficients. If one such term is $x^{a}y^{b}z^{c}$, then b and c are both even or both odd. $a + b + c <= 2006$, so we can do casework on a:$a = 0, b = 0, 2, 4... 2006, c = 2006, 2004...0$ 1004 possibilities.$a = 1, b = 1, 3, 5... 2005$ 1003$a = 2, b = 0, 2, 4... 2004$ 1003$a = 3, b = 1, 3, 5... 2003$ 1002$a = 4, b = 0, 2, 4... 2002$ 1002$a = 5, b = 1, 3, 5... 2001$ 1001.................................$a = 2004, b = 0, 2$ 2$a = 2005, b = 1$ 1$a = 2006, b = 0$ 1.So there are $2(1+2+3...+1003) + 1004$ total, which is $1003*1004 + 1004$ = $1004^2$ = 1008016 $\implies \boxed{D}$




\textbf{Solution 7(Casu Martzu)}

As in Solution 1, note that if we assume no terms are cancelled, by stars and bars there will be $\dbinom{2006+2}{3}=2015028.$ Now, to get a lower bound, note that at most half of the terms will get their sign changed, as there is about a 50-50 chance that we will get an odd number of $-y$ and $-z$ terms. This eliminates options $A,B,E.$ Finally, since there are terms where we have an odd number of both $y$ and $z$ terms, the chances are slightly higher than 50-50, leaving us with $D$ as our only option.

 


	\item (\textit{AMC 12 12A 2006}) How many non-empty subsets $S$ of $\lbrace 1,2,3,\ldots ,15\rbrace$ have the following two properties? 


 $(1)$  No two consecutive integers belong to $S$.


 $(2)$  If $S$ contains $k$ elements, then $S$ contains no number less than $k$.


 $\mathrm{(A) \ } 277\qquad \mathrm{(B) \ } 311\qquad \mathrm{(C) \ } 376\qquad \mathrm{(D) \ } 377\qquad \mathrm{(E) \ }  405$

 \textbf{Answer: }E



\textbf{Solution 1}

This question can be solved fairly directly by casework and pattern-finding.  We give a somewhat more general attack, based on the solution to the following problem:

 How many ways are there to choose $k$ elements from an ordered $n$ element \href{https://artofproblemsolving.com/wiki/index.php?title=Set}{set} without choosing two consecutive members?

 You want to choose $k$ numbers out of $n$ with no consecutive numbers. For each configuration, we can subtract $i-1$ from the $i$-th element in your subset. This converts your configuration into a configuration with $k$ elements where the largest possible element is $n-k+1$, with no restriction on consecutive numbers. Since this process is easily reversible, we have a \href{https://artofproblemsolving.com/wiki/index.php?title=Bijection}{bijection}.Without consideration of the second condition, we have:${15 \choose 1} + {14 \choose 2} + {13 \choose 3} + ... + {9 \choose 7} + {8 \choose 8}$

 Now we examine the second condition. It simply states that no element in our original configuration (and hence also the modified configuration, since we don't move the smallest element) can be less than $k$, which translates to subtracting $k - 1$ from the "top" of each \href{https://artofproblemsolving.com/wiki/index.php?title=Binomial\_coefficient}{binomial coefficient}.Now we have, after we cancel all the terms ${n \choose k}$ where $n < k$, ${15 \choose 1} + {13 \choose 2} + {11 \choose 3} + {9 \choose 4} + {7 \choose 5}=  15 + 78 + 165 + 126 + 21 = \boxed{405} \Longrightarrow \mathrm{(E)}$




\textbf{Solution 2}

Another way of visualizing the solution above would be to use $|$'s and $-$'s. Denote $|$ as the numbers we have chosen and $-$ as other numbers. Taking an example, assuming we are picking two numbers, we imagine the shape $| - |$. This notation forces a number between the two chosen numbers, which blocks the two numbers we picked from being consecutive. Now we consider the orientations with this shape. We have $15 - 3 = 12$ remaining numbers.

 We need to find the number of ways to place the remaining $12 -$'s. We can find this by utilizing stars and bars, with the following marker being placed to represent groups: *| - *|*. Now, we have to place $12$ numbers within $3$ groups, which is ${14 \choose 2}$. The same concept can be used for the remaining numbers. The rest of the solution continues as above.

 Solution by: Everyoneintexas

 



\textbf{Solution 3}

We have the same setup as in the previous solution.

 Note that if $n < 2k - 1$, the answer will be 0.  Otherwise, the $k$ elements we choose define $k + 1$ boxes (which divide the nonconsecutive numbers) into which we can drop the $n - k$ remaining elements, with the caveat that each of the middle $k - 1$ boxes must have at least one element (since the numbers are nonconsecutive).  This is equivalent to dropping $n - 2k + 1$ elements into $k + 1$ boxes, where each box is allowed to be empty.  And this is equivalent to arranging $n - k + 1$ objects, $k$ of which are dividers, which we can do in $F(n, k) = {n - k + 1 \choose k}$ ways.

 Now, looking at our original question, we see that the thing we want to calculate is just $F(15, 1) + F(14, 2) + F(13, 3) + F(12, 4) + F(11, 5) = {15\choose 1} + {13\choose2} + {11\choose 3} + {9\choose 4} + {7 \choose 5} = 15 + 78 + 165 + 126 + 21  = 405 \Longrightarrow \mathrm{(E)}$




\textbf{Solution 4}

Let $s$ be the numbers of elements in a subset. First we examine the second condition. No elements less than $s$ can be put in a subset of size $s$, therefore the "lowest" element that can go into the subset is $s$, whereas the "highest" element that can go into the subset is $15$. This is a total of $15-s+1$ or $16-s$ possible elements. 

 Now we consider the first condition. No consecutive elements are allowed. This means that if an element $e$ is put into the subset, both $e-1$ and $e+1$ are no longer possibilities. Assume that all subsets are ordered from least to greatest (we are looking for the number of combinations, so we can order these combinations however we want). Then the first element will be $s$ (as a reminder, the lowest possible element in a subset is $s$), the second element will be at least $s+2$, and so on. After $s$ elements are chosen, we will have skipped $s-1$ elements (these are the elements that were "eliminated" as they were consecutive). Therefore, we ignore exactly $s-1$ elements (if we ignore more or less elements, then $s$ changes) Since we must ignore $s-1$ elements, we can simply remove those beforehand. ($16-s)-(s-1) = 17-2s$ possible elements

 Now we look for the bounds of $s$. We are looking for non-empty subsets, so $s\ge1$. If $s$ is too large, there will not be enough non-consecutive terms between $s$ and $15$. Specifically, the highest element in a subset using "optimal" selection will be $s+2(s-1)$ or $3s-2$. If $3s-1>15$, that means s is too large. Therefore $3s-2\le15$; solving for $s$ yields $s\le5$. Now we know that $1 \le s \le 5$.

 We want to know the number of ways to arrange $17-2s$ "balls" into $s$ identical "boxes" with at most $1$ ball per box, for $s=1$ to $s=5$. This is equivalent to $\sum_{s=1}^{5}{17-2s \choose s} = 405$, or $\boxed{\text{E}}$.

 



\textbf{Solution 5 (Possibly Clever bash)}

We will split the problem into cases, and maybe one could then generalize this to arbitrary $n$. 

 $\bold { Case 1}:$ ($n=15, k=1$). Then this is easy. We have $\binom{15}{1}=15$.

 $\bold{Case 2}:$ ($n=14, k=2$). Now we have something tricky. To get a good grasp on this case, let us consider the smallest element; $2$, in the first spot. We then have $15-1=14$ numbers left. However, we cannot have the digit $3$. Hence we have to choose $2$ numbers from $14-1=13$ integers left. This can be done in $\binom{13}{2}=78$ ways.

 $\bold{Case 3}:$ ($n=13, k=3$). As in the last case, we can use the smallest element to get a good grasp on this case. Put the digit $3$ in the first spot. Then there are $11$ integers left for the second spot. But in total, we have to avoid $2$ elements for each subset of cardinality $3$ ${a,b,c}$. So we have to choose $3$ elements from $13-2=11$ elements, which can be done in $\binom{11}{3}=165$ ways. 

 $\bold{Case 4}:$ ($n=12, k=4$). As in the previous case, we have to avoid $3$ numbers, and choose $4$ from them. Which is choosing $4$ elements from $12-3=9$ elements, which can be done in $\binom{9}{4}=126$ ways.

 $\bold{Case 5}:$ ($n=11, k=4$). Now you are accustomed to the strategy. We remove $4$ integers leaving $\binom{7}{5}=21$ ways.

 Adding up all of the cases yields $405$, as in $\boxed{E)405}$. SHEN KISLAY KAI~~2023

 Note also that we only went up to cardinality or $k=5$, or else for greater $k$, we would always have consecutive numbers.

 


\end{enumerate}